\documentclass[a4paper,11pt]{article}
\usepackage[utf8]{inputenc}
\usepackage[T1]{fontenc}
\usepackage[french]{babel}
\usepackage[top=1.5cm, bottom=1.cm, left=1.5cm, right=1.5cm]{geometry}
\usepackage{array}
\usepackage{tabularx}
\usepackage{colortbl}
\usepackage{xcolor}
\usepackage{datatool}
\usepackage{ifthen}

\pagestyle{empty}

% --- SOLUTION ALTERNATIVE POUR LE SÉPARATEUR ---
% On définit le séparateur directement dans la commande de chargement
\DTLloaddb[noheader=false, separator=;]{eleves}{g1.csv}

\def\mycourse{CIEL 1}
\def\mydate{04 mai 2026}
\def\mytime{14H00}
\def\mygroup{1}
\def\myroom{29}
\def\myprof{RICART THOMAS}

\begin{document}

\begin{center}
    {\large \textbf{Contrôle en Cours de Formation -- Situation en Centre de Formation}}\\[0.4cm]
    {\LARGE \textbf{CONVOCATION}}
\end{center}

\vspace{0.3cm}

\noindent
\textbf{Matière :} MATHÉMATIQUES \hfill \textbf{Épreuve :} CCF situation 1 

\vspace{0.3cm}
\noindent
\textbf{Classe :} BTS \mycourse -- GROUPE \mygroup 

\vspace{0.5cm}
\noindent
\textbf{Nom du professeur faisant passer l'épreuve :} \myprof 

\vspace{0.2cm}
\noindent
\textbf{Lieu de l'épreuve :} SALLE \myroom 

\vspace{0.6cm}

\renewcommand{\arraystretch}{1.8}
\begin{tabularx}{\textwidth}{|c|X|c|c|p{3.5cm}|p{3.5cm}|}
\hline
\rowcolor{gray!25}
\textbf{N°} & \textbf{Nom Prénom ÉTUDIANT} & \textbf{DATE} & \textbf{Horaires} & \textbf{Pris connaissance le} & \textbf{Émargement ÉLÈVE} \\
\hline
% Boucle ultra-propre : on n'affiche la ligne QUE si \nom n'est pas vide
\DTLforeach{eleves}{\nom=Eleve}{%
    \ifthenelse{\equal{\nom}{}}{}{%
        \DTLcurrentindex & \nom & \mydate & \mytime & & \tabularnewline\hline
    }%
}%
\end{tabularx}

\vspace{1cm}
\noindent
\textbf{Document transmis le} \underline{\hspace{4cm}} 

\vspace{1.5cm}
\noindent
\begin{minipage}[t]{0.45\textwidth}
    \textbf{Le Professeur,}\\[2.5cm] 
\end{minipage}
\hfill
\begin{minipage}[t]{0.45\textwidth}
    \textbf{Le DDF,}\\[2.5cm] 
\end{minipage}

\end{document}